\documentclass[12pt,a4paper]{article}
\usepackage[utf8]{inputenc}
\usepackage[T1]{fontenc}
\usepackage[spanish]{babel}
\usepackage{geometry}
\usepackage{xcolor}
\usepackage{listings}
\usepackage{booktabs}
\usepackage{hyperref}
\usepackage{fancyhdr}
\usepackage{titlesec}
\usepackage{array}
\usepackage{longtable}
\usepackage{parskip}
\usepackage{mdframed}
\usepackage{enumitem}
\usepackage{tikz}
\usetikzlibrary{arrows.meta, shapes.geometric, positioning}

\geometry{top=2.5cm, bottom=2.5cm, left=3cm, right=3cm}

\definecolor{codeblue}{RGB}{41, 98, 255}
\definecolor{codegray}{RGB}{128, 128, 128}
\definecolor{codegreen}{RGB}{38, 166, 65}
\definecolor{bgcode}{RGB}{245, 247, 250}
\definecolor{accent}{RGB}{0, 112, 243}
\definecolor{headerblue}{RGB}{15, 30, 70}
\definecolor{warnyellow}{RGB}{200, 140, 0}

\lstset{
  backgroundcolor=\color{bgcode},
  basicstyle=\ttfamily\footnotesize,
  breaklines=true,
  captionpos=b,
  commentstyle=\color{codegray},
  keywordstyle=\color{codeblue}\bfseries,
  stringstyle=\color{codegreen},
  numberstyle=\tiny\color{codegray},
  numbers=left,
  numbersep=8pt,
  frame=single,
  framesep=4pt,
  rulecolor=\color{codegray!40},
  tabsize=2,
  showstringspaces=false,
  xleftmargin=12pt,
}

\titleformat{\section}{\Large\bfseries\color{headerblue}}{}{0em}{\thesection\quad}[\titlerule]
\titleformat{\subsection}{\large\bfseries\color{accent}}{}{0em}{\thesubsection\quad}
\titleformat{\subsubsection}{\normalsize\bfseries}{}{0em}{\thesubsubsection\quad}

\pagestyle{fancy}
\fancyhf{}
\fancyhead[L]{\small\textcolor{codegray}{SaukCode.cl --- Documentación Técnica}}
\fancyhead[R]{\small\textcolor{codegray}{Vol. III: Utilidades y Datos}}
\fancyfoot[C]{\small\textcolor{codegray}{\thepage}}
\renewcommand{\headrulewidth}{0.4pt}

\mdfdefinestyle{infobox}{
  linecolor=accent, linewidth=1.5pt, backgroundcolor=accent!5,
  innertopmargin=8pt, innerbottommargin=8pt,
  innerleftmargin=12pt, innerrightmargin=12pt, roundcorner=4pt,
}

\mdfdefinestyle{warnbox}{
  linecolor=warnyellow, linewidth=1pt, backgroundcolor=warnyellow!8,
  innertopmargin=8pt, innerbottommargin=8pt,
  innerleftmargin=10pt, innerrightmargin=10pt, roundcorner=3pt,
}

\hypersetup{colorlinks=true, linkcolor=accent, urlcolor=codeblue}

\begin{document}

% Portada
\thispagestyle{empty}
\begin{center}
  \vspace*{3cm}
  {\Huge\bfseries\color{headerblue} SaukCode.cl}\\[0.5cm]
  {\Large\color{accent} Documentación Técnica}\\[0.3cm]
  {\large\color{codegray} Volumen III: Utilidades, Hooks y Capa de Datos}\\[2cm]
  \rule{10cm}{1.5pt}\\[0.5cm]
  {\normalsize Custom Hooks \textbullet\ Utils \textbullet\ Data Layer \textbullet\ Content Structure}\\[0.5cm]
  \rule{10cm}{1.5pt}\\[2cm]
  {\small\color{codegray} Generado: Mayo 2026}
\end{center}
\newpage

\tableofcontents
\newpage

% =====================================================
\section{Custom Hooks}
% =====================================================

Los tres hooks personalizados implementan la lógica reutilizable de las tablas interactivas del sitio. Siguen el patrón de composición: cada hook hace una sola cosa y se combina con los demás.

\subsection{usePagination}

\textbf{Ruta}: \texttt{src/hooks/usePagination.jsx}

Gestiona el estado y cálculos de paginación para cualquier array de datos.

\subsubsection{API}

\begin{table}[h]
\centering
\begin{tabular}{lll}
\toprule
\textbf{Input} & \textbf{Tipo} & \textbf{Descripción} \\
\midrule
\texttt{data} & \texttt{Array} & Datos a paginar \\
\texttt{itemsPerPage} & \texttt{number} & Elementos por página (default: 10) \\
\texttt{resetTriggers} & \texttt{Array} & Valores que resetean a página 1 \\
\bottomrule
\end{tabular}
\end{table}

\begin{table}[h]
\centering
\begin{tabular}{lll}
\toprule
\textbf{Output} & \textbf{Tipo} & \textbf{Descripción} \\
\midrule
\texttt{currentPage} & \texttt{number} & Página actual (validada) \\
\texttt{totalPages} & \texttt{number} & Total de páginas \\
\texttt{totalItems} & \texttt{number} & Total de elementos \\
\texttt{paginatedData} & \texttt{Array} & Slice de datos para la página actual \\
\texttt{hasNextPage} & \texttt{boolean} & ¿Hay página siguiente? \\
\texttt{hasPreviousPage} & \texttt{boolean} & ¿Hay página anterior? \\
\texttt{isFirstPage} & \texttt{boolean} & ¿Es la primera página? \\
\texttt{isLastPage} & \texttt{boolean} & ¿Es la última página? \\
\texttt{goToPage(n)} & \texttt{function} & Navega a la página n (validado) \\
\texttt{goToFirstPage()} & \texttt{function} & Navega a la primera página \\
\texttt{goToLastPage()} & \texttt{function} & Navega a la última página \\
\texttt{goToNextPage()} & \texttt{function} & Navega a la siguiente página \\
\texttt{goToPreviousPage()} & \texttt{function} & Navega a la anterior \\
\texttt{getPageNumbers(max)} & \texttt{function} & Array de números de página visibles \\
\bottomrule
\end{tabular}
\end{table}

\subsubsection{Detalles de Implementación}

\begin{itemize}
  \item Usa \texttt{useMemo} para calcular toda la información de paginación en una sola pasada
  \item Los \texttt{resetTriggers} se serializan con \texttt{JSON.stringify} para comparación estable en el \texttt{useEffect}
  \item La función \texttt{getPageNumbers} implementa una ventana deslizante: si hay más páginas que \texttt{maxVisible}, muestra las páginas alrededor de la actual
  \item Todos los valores de página están \textbf{validados con clamp} (\texttt{Math.max(1, Math.min(...))}) para evitar páginas fuera de rango
\end{itemize}

\begin{lstlisting}[language=JavaScript, caption={Uso típico en un componente de tabla}]
const { paginatedData, currentPage, totalPages,
        goToPage, getPageNumbers } = usePagination({
  data: filteredData,   // Ya filtrado por useFilters
  itemsPerPage: 20,
  resetTriggers: [searchTerm, difficultyFilter] // Reset al filtrar
});
\end{lstlisting}

\subsection{useFilters}

\textbf{Ruta}: \texttt{src/hooks/useFilters.jsx}

Gestiona múltiples filtros simultáneos sobre un dataset. Diseñado principalmente para los datos de LeetCode pero reutilizable para cualquier colección de objetos con campos \texttt{title}, \texttt{topics}, \texttt{difficulty} y \texttt{categories}.

\subsubsection{Filtros Disponibles}

\begin{table}[h]
\centering
\begin{tabular}{lll}
\toprule
\textbf{Filtro} & \textbf{Estado} & \textbf{Comportamiento} \\
\midrule
Búsqueda (\texttt{searchTerm}) & \texttt{string} & Busca en title, id y topics \\
Tópico (\texttt{topicFilter}) & \texttt{string} & Exacto: item debe incluir el tópico \\
Dificultad (\texttt{difficultyFilter}) & \texttt{string} & Exacto: Easy/Medium/Hard \\
Categoría (\texttt{categoryFilter}) & \texttt{string} & Exacto: item debe incluir la categoría \\
\bottomrule
\end{tabular}
\end{table}

\subsubsection{Normalización de Texto}

El hook normaliza el texto para búsqueda sin diferencias de tildes:

\begin{lstlisting}[language=JavaScript]
const normalizeText = (text) => {
  return text
    .toLowerCase()
    .normalize('NFD')
    .replace(/[̀-ͯ]/g, '');
  // "Árbol" -> "arbol"
};
\end{lstlisting}

\subsubsection{Lógica de Búsqueda Multi-palabra}

Si el \texttt{searchTerm} contiene más de una palabra, cada palabra debe estar presente en al menos uno de los tópicos:

\begin{lstlisting}[language=JavaScript]
const searchWords = normalizedSearchTerm
  .split(/\s+/).filter(w => w.length > 0);
if (searchWords.length > 1) {
  return searchWords.every(word => normalizedTopic.includes(word));
}
\end{lstlisting}

\subsubsection{Mapeo de Dificultades}

El hook incluye mapeos bidireccionales para normalizar las dificultades entre el formato interno (EN) y el mostrado al usuario (ES):

\begin{lstlisting}[language=JavaScript]
const difficultyMap = {
  'Easy': 'Fácil', 'Medium': 'Medio', 'Hard': 'Difícil'
};
const reverseDifficultyMap = {
  'Fácil': 'Easy', 'Medio': 'Medium', 'Difícil': 'Hard'
};
\end{lstlisting}

\subsubsection{Salidas Computadas}

Además de los datos filtrados, el hook extrae automáticamente todos los valores únicos disponibles para poblar los selects de filtro:
\begin{itemize}
  \item \texttt{allTopics}: todos los tópicos únicos del dataset (ordenados)
  \item \texttt{allCategories}: todas las categorías únicas del dataset (ordenadas)
\end{itemize}

\subsection{useTableSort}

\textbf{Ruta}: \texttt{src/hooks/useTableSort.jsx}

Gestiona el estado de ordenamiento de columnas para tablas de datos.

\subsubsection{API}

\begin{lstlisting}[language=JavaScript]
const { sortConfig, sortedData, requestSort } = useTableSort(
  data,              // Array de [id, object] (formato de Object.entries)
  { key: 'id', direction: 'asc' } // Ordenamiento por defecto
);
\end{lstlisting}

\subsubsection{Columnas de Ordenamiento}

Actualmente implementa ordenamiento para cuatro columnas de LeetCode:

\begin{table}[h]
\centering
\begin{tabular}{ll}
\toprule
\textbf{Clave} & \textbf{Criterio de ordenamiento} \\
\midrule
\texttt{id} & Numérico: extrae dígitos del ID con \texttt{/\textbackslash D/g} \\
\texttt{title} & Lexicográfico con \texttt{localeCompare} \\
\texttt{topics} & Lexicográfico por el primer tópico del array \\
\texttt{difficulty} & Por peso: Easy=1, Medium=2, Hard=3 \\
\bottomrule
\end{tabular}
\end{table}

\subsubsection{Toggle de Dirección}

\begin{lstlisting}[language=JavaScript]
const requestSort = (key) => {
  let direction = 'asc';
  if (sortConfig.key === key && sortConfig.direction === 'asc') {
    direction = 'desc';  // Segundo clic invierte
  }
  setSortConfig({ key, direction });
};
\end{lstlisting}

\subsubsection{Composición de los Tres Hooks}

\begin{lstlisting}[language=JavaScript, caption={Patrón de composición típico en ProblemTable.jsx}]
// 1. Filtrar datos
const { filteredData, searchTerm, setSearchTerm, ... } =
  useFilters(problems);

// 2. Ordenar datos filtrados
const { sortedData, sortConfig, requestSort } =
  useTableSort(filteredData);

// 3. Paginar datos ordenados+filtrados
const { paginatedData, currentPage, totalPages, goToPage } =
  usePagination({
    data: sortedData,
    itemsPerPage: 20,
    resetTriggers: [searchTerm, difficultyFilter]
  });
\end{lstlisting}

% =====================================================
\section{Utilidades (src/utils/)}
% =====================================================

\subsection{contentUtils.js}

\textbf{Propósito}: Interfaz entre el filesystem de contenido y el sistema de rutas de Next.js.

\subsubsection{extractSlugsFromCourseData(courseData)}

Extrae todos los slugs disponibles de un objeto de datos de curso recorriendo recursivamente sus arrays.

\begin{mdframed}[style=infobox]
\textbf{Estructuras soportadas:}
\begin{itemize}[noitemsep]
  \item \texttt{courseData.classData[].description[].href}
  \item \texttt{courseData.lectureData[].description[].href}
  \item \texttt{courseData.evaluationsData[].description.href}
  \item \texttt{courseData.modules[].classes[].link}
  \item \texttt{courseData.units[].sections[].link}
\end{itemize}
\end{mdframed}

\subsubsection{getLessonSlugs(section, platform, courseId)}

\begin{lstlisting}[language=JavaScript]
// Ejemplo:
getLessonSlugs('inacap', null, 'ti3v53')
// -> ['lectura01', 'lectura02', ..., 'lectura10']
\end{lstlisting}

Lee directamente el filesystem para obtener los slugs disponibles (más confiable que parsear datos). Filtra archivos \texttt{.mdx} y elimina la extensión.

\subsubsection{getAllLessonPaths(section, platforms)}

Genera todos los pares \texttt{\{courseId, slug\}} o \texttt{\{platform, courseId, slug\}} para \texttt{generateStaticParams()}:

\begin{lstlisting}[language=JavaScript]
// Sin plataformas (INACAP, idiomas)
getAllLessonPaths('inacap')
// -> [{courseId:'ti3v11', slug:'lectura01'}, ...]

// Con plataformas (bootcamp, elearning)
getAllLessonPaths('bootcamp', ['codfacilito', 'udd'])
// -> [{platform:'codfacilito', courseId:'cft001', slug:'clase01'}, ...]
\end{lstlisting}

\subsubsection{lessonExists(section, platform, courseId, slug)}

Validación de existencia de archivo MDX. Usa \texttt{fs.existsSync}. Útil para evitar errores 404 en links.

\subsubsection{getMDXImportPath(section, platform, courseId, slug)}

Retorna la ruta de import dinámica para \texttt{next/dynamic} o el import directo:

\begin{lstlisting}[language=JavaScript]
getMDXImportPath('inacap', null, 'ti3v53', 'lectura10')
// -> '@/content/inacap/ti3v53/lectura10.mdx'
\end{lstlisting}

\subsection{linkBuilder.js}

\textbf{Propósito}: Centralizar la construcción de URLs para evitar inconsistencias entre componentes.

\begin{table}[h]
\centering
\begin{tabular}{ll}
\toprule
\textbf{Función} & \textbf{Ejemplo de salida} \\
\midrule
\texttt{buildLessonLink(s, c, slug)} & \texttt{/elearning/cft002/clase01} \\
\texttt{buildLanguageLessonLink(lang, href)} & \texttt{/chino/ch001/lesson11} \\
\texttt{buildCourseLink(section, courseId)} & \texttt{/bootcamp/cft001} \\
\texttt{buildCourseResourceLink(s, c, path)} & \texttt{/inacap/ti3v11/docs/resumen.pdf} \\
\bottomrule
\end{tabular}
\end{table}

Todas las funciones incluyen validación de parámetros y retornan \texttt{'\#'} con un \texttt{console.warn} si faltan argumentos requeridos.

\begin{mdframed}[style=warnbox]
\textbf{Nota de diseño:} \texttt{buildLanguageLessonLink} normaliza el \texttt{href} eliminando el slash inicial si existe, para evitar doble slash: \texttt{/chino//ch001/lesson} → \texttt{/chino/ch001/lesson}.
\end{mdframed}

\subsection{animations.js}

\textbf{Propósito}: Centralizar todas las configuraciones de Framer Motion para mantener consistencia visual.

\subsubsection{Variantes Base}

\begin{table}[h]
\centering
\begin{tabular}{ll}
\toprule
\textbf{Export} & \textbf{Efecto} \\
\midrule
\texttt{fadeIn} & Aparición con opacidad \\
\texttt{fadeInUp} & Aparición desde abajo (y: 20→0) \\
\texttt{fadeInDown} & Aparición desde arriba (y: -20→0) \\
\texttt{fadeInLeft} & Aparición desde la izquierda \\
\texttt{fadeInRight} & Aparición desde la derecha \\
\texttt{scaleIn} & Zoom in (scale: 0.9→1) \\
\texttt{scrollRevealBlur} & Aparición con desenfoque (blur: 10px→0) \\
\texttt{scrollRevealScale} & Aparición con escala (scale: 0.85→1) \\
\texttt{pageTransition} & Transición estándar de página \\
\texttt{menuSlide} & Panel lateral (x: 100\%→0, spring) \\
\texttt{menuOverlay} & Fondo oscuro del menú (fadeIn) \\
\texttt{menuItemVariants} & Items del menú con stagger \\
\texttt{parallaxSlow} & Parallax suave de elementos de fondo \\
\bottomrule
\end{tabular}
\end{table}

\subsubsection{Factories}

\begin{lstlisting}[language=JavaScript]
// Variante personalizada de fade in
const myVariant = createFadeIn('left', 40);

// Contenedor stagger personalizado
const myContainer = createStaggerContainer({
  stagger: 0.15, delay: 0.2, when: 'beforeChildren'
});

// Contenedor stagger estándar
const container = staggerContainer(0.1, 0); // stagger=0.1, delay=0
\end{lstlisting}

\subsubsection{Configuraciones de Viewport}

\begin{lstlisting}[language=JavaScript]
// Animación de scroll estándar (20% visible, -50px margen)
<motion.div
  initial="initial"
  whileInView="animate"
  viewport={viewportConfig}
  variants={fadeInUp}
>

// Trigger temprano (10% visible, -100px margen)
<motion.div viewport={viewportEarly} ... >
\end{lstlisting}

\subsubsection{Transiciones Predefinidas}

\begin{lstlisting}[language=JavaScript]
export const transitions = {
  spring:  { type: 'spring', damping: 20, stiffness: 300 },
  smooth:  { duration: 0.4, ease: [0.25, 0.1, 0.25, 1] },
  fast:    { duration: 0.2, ease: 'easeOut' },
  slow:    { duration: 0.6, ease: [0.25, 0.1, 0.25, 1] },
};
\end{lstlisting}

La curva \texttt{[0.25, 0.1, 0.25, 1]} es una curva \textit{ease-in-out cubic bezier} que produce animaciones naturales y modernas.

\subsection{Otras Utilidades}

\subsubsection{bootcampUtils.js}

Funciones para acceder a datos de cursos de bootcamp:
\begin{itemize}
  \item \texttt{getCourseData(platform, courseId)}: Retorna el objeto de datos del curso
  \item \texttt{extractBootcampSlugs(courseData)}: Extrae slugs de la estructura de módulos/clases
  \item \texttt{buildBootcampPaths()}: Genera todos los paths para \texttt{generateStaticParams}
\end{itemize}

\subsubsection{elearningUtils.js}

Agregación de datos de plataformas e-learning:
\begin{itemize}
  \item \texttt{getPlatformData(platform)}: Metadatos de una plataforma
  \item \texttt{getAllPlatforms()}: Lista todas las plataformas disponibles
  \item \texttt{getCoursesByPlatform(platform)}: Cursos de una plataforma específica
\end{itemize}

\subsubsection{idiomasUtils.js}

Helpers para cursos de idiomas:
\begin{itemize}
  \item \texttt{getLanguageData(language)}: Datos de un idioma (cursos, módulos)
  \item \texttt{getLanguageLessons(language, courseId)}: Lecciones de un módulo
\end{itemize}

\subsubsection{courseUtils.js}

Helpers específicos de INACAP:
\begin{itemize}
  \item \texttt{getCourseById(courseId)}: Datos de una asignatura por código
  \item \texttt{getCoursesByCategory(category)}: Asignaturas filtradas por categoría
  \item \texttt{getSemesterCourses(semester)}: Asignaturas de un semestre
\end{itemize}

\subsubsection{leetcodeUtils.js}

Filtrado y acceso a problemas:
\begin{itemize}
  \item \texttt{getProblemsArray()}: Convierte el objeto de problemas a array
  \item \texttt{filterByDifficulty(problems, difficulty)}: Filtra por nivel
  \item \texttt{filterByTopic(problems, topic)}: Filtra por tópico
\end{itemize}

% =====================================================
\section{Capa de Datos (src/data/)}
% =====================================================

\subsection{Patrón de Estructura de Datos de Curso}

Todos los cursos (independientemente de la sección) siguen una estructura base común con variaciones según la plataforma:

\subsubsection{Estructura INACAP}

\begin{lstlisting}[language=JavaScript, caption={Patrón de datos INACAP (ti3v53.js)}]
export const ti3v53Data = {
  code: 'TI3V53',
  title: 'Nombre de la Asignatura',
  semester: 5,
  category: 'programacion',

  classData: [
    {
      title: 'Clases',
      description: [
        { name: 'Lectura 01: Introducción', href: '/lectura01' },
        { name: 'Lectura 02: Conceptos', href: '/lectura02', done: true },
      ]
    }
  ],

  evaluationsData: [
    {
      title: 'Evaluación 1',
      description: { name: 'Prueba Diagnóstica', href: '/diagnostica' },
      date: '2025-03-15',
      weight: 0.2,
    }
  ],

  lectureData: [
    {
      title: 'Material Adicional',
      description: [
        { name: 'Guía de Ejercicios', href: '/guia01' }
      ]
    }
  ]
};
\end{lstlisting}

\subsubsection{Estructura Bootcamp (CódigoFacilito)}

\begin{lstlisting}[language=JavaScript, caption={Patrón de datos Bootcamp (cft001.js)}]
export const cft001Data = {
  id: 'cft001',
  title: 'Nombre del Bootcamp',
  platform: 'codfacilito',

  modules: [
    {
      name: 'Módulo 1: Fundamentos',
      classes: [
        {
          classNumber: '01',
          name: 'Introducción al curso',
          link: 'clase01',
          done: true,
          exercises: { link: 'ejercicio01' }
        }
      ]
    }
  ]
};
\end{lstlisting}

\subsubsection{Estructura E-Learning}

\begin{lstlisting}[language=JavaScript, caption={Patrón de datos E-Learning}]
export const courseData = {
  id: 'udemy001',
  title: 'Nombre del Curso',
  platform: 'udemy',
  instructor: 'Nombre Instructor',
  rating: 4.8,

  units: [
    {
      name: 'Sección 1',
      sections: [
        { name: 'Clase 01', link: 'clase01', done: true }
      ]
    }
  ]
};
\end{lstlisting}

\subsubsection{Estructura LeetCode}

\begin{lstlisting}[language=JavaScript, caption={Estructura problems.js (fragmento)}]
export const problems = {
  'lc001': {
    title: 'Two Sum',
    difficulty: 'Easy',
    topics: ['Array', 'Hash Table'],
    categories: ['Fundamentals'],
    done: true,
    link: 'two-sum'
  },
  'lc002': {
    title: 'Add Two Numbers',
    difficulty: 'Medium',
    topics: ['Linked List', 'Recursion'],
    categories: ['Linked List'],
    done: false,
    link: 'add-two-numbers'
  }
};
\end{lstlisting}

\subsection{homeData.js}

Datos de la página de inicio. Define las tres secciones visibles en el \texttt{CardCarousel}:

\begin{lstlisting}[language=JavaScript]
export const sections = [
  {
    id: 'programacion',
    title: 'Programación',
    description: '...',
    color: '#0070f3',
    link: '/inacap',
    subsections: ['INACAP', 'Bootcamp', 'E-Learning']
  },
  {
    id: 'idiomas',
    title: 'Idiomas',
    description: '...',
    color: '#7928ca',
    link: '/ingles',
    subsections: ['Inglés', 'Alemán', 'Coreano', ...]
  },
  { id: 'sobremi', title: 'Sobre Mí', link: '/portfolio', ... }
];
\end{lstlisting}

\subsection{calendar25.js}

Calendario académico INACAP 2025. Array de eventos con:

\begin{lstlisting}[language=JavaScript]
{ date: '2025-03-10', title: 'Inicio de clases', type: 'start' }
{ date: '2025-04-18', title: 'Feriado de Semana Santa', type: 'holiday' }
{ date: '2025-06-20', title: 'Prueba Parcial 1', type: 'exam' }
\end{lstlisting}

Tipos de evento: \texttt{start}, \texttt{holiday}, \texttt{exam}, \texttt{project}, \texttt{end}.

% =====================================================
\section{Estructura de Contenido MDX}
% =====================================================

\subsection{Organización del Filesystem}

\begin{lstlisting}[language=bash]
src/content/
├── inacap/
│   ├── ti3v11/           # Una carpeta por asignatura
│   │   ├── lectura01.mdx
│   │   └── lectura02.mdx
│   └── ti3v53/
│       ├── lectura01.mdx
│       └── diagnostica.mdx
├── bootcamp/
│   ├── codfacilito/      # Subcarpeta por plataforma
│   │   └── cft001/
│   │       ├── clase01.mdx
│   │       └── ejercicio01.mdx
│   └── udd/
│       └── udd001/
├── elearning/
│   ├── udemy/
│   ├── devtalles/
│   ├── codfacilito/
│   ├── platzi/
│   └── coursera/
├── idiomas/
│   ├── ingles/
│   ├── aleman/
│   ├── coreano/
│   ├── chino/
│   ├── japones/
│   └── ruso/
└── leetcode/             # Sin subcarpetas
    ├── two-sum.mdx
    └── add-two-numbers.mdx
\end{lstlisting}

\subsection{Anatomía de un Archivo MDX}

\begin{lstlisting}[language=HTML, caption={Estructura típica de un archivo MDX de lectura}]
# Título de la Lectura

Introducción en párrafo con **negrita** y *cursiva*.

## Sección Principal

Contenido de la sección con enlace interno:
<SmartLink href="/inacap/ti3v53/lectura09">Lectura anterior</SmartLink>

### Subsección con Código

<CodeBlock language="python" filename="ejemplo.py">
def bubble_sort(arr):
    n = len(arr)
    for i in range(n):
        for j in range(n - i - 1):
            if arr[j] > arr[j + 1]:
                arr[j], arr[j + 1] = arr[j + 1], arr[j]
</CodeBlock>

## Fórmulas Matemáticas

La complejidad es $O(n^2)$ en el peor caso.

$$\sum_{i=1}^{n} i = \frac{n(n+1)}{2}$$

## Tabla de Datos

<SmartTable
  columns={[
    { key: 'caso', label: 'Caso' },
    { key: 'complejidad', label: 'Complejidad', formula: true }
  ]}
  data={[
    { caso: 'Mejor', complejidad: 'O(n)' },
    { caso: 'Promedio', complejidad: 'O(n^2)' }
  ]}
/>

## Evaluación

<Quiz>
  <QuizQuestion
    question="¿Cuál es la complejidad del Bubble Sort?"
    options={['O(n)', 'O(n log n)', 'O(n²)', 'O(1)']}
    answer={2}
    explanation="Bubble Sort compara cada par adyacente repetidamente"
  />
</Quiz>

## Solución del Ejercicio

<Solution>
  El algoritmo de Bubble Sort usa $O(1)$ espacio adicional.
</Solution>
\end{lstlisting}

% =====================================================
\section{Flujo Completo de Datos}
% =====================================================

El siguiente diagrama muestra el flujo completo desde los archivos de datos hasta la página renderizada:

\begin{lstlisting}[language=bash, caption={Flujo de datos en build time}]
BUILD TIME:
  data/inacap/ti3v53.js
      └─> utils/contentUtils.js::getAllLessonPaths('inacap')
              └─> filesystem scan: content/inacap/ti3v53/*.mdx
                      └─> [{courseId:'ti3v53', slug:'lectura10'}, ...]
                              └─> generateStaticParams()
                                      └─> Pre-genera /inacap/ti3v53/lectura10

RUNTIME (página pre-generada):
  /inacap/ti3v53/lectura10
      └─> import('@/content/inacap/ti3v53/lectura10.mdx')
              └─> <CustomMDXProvider>
                      └─> <Content /> (MDX compilado)
                              └─> SmartTable, CodeBlock, Quiz...
\end{lstlisting}

\subsection{Carga Dinámica de MDX en Runtime}

Las páginas de lección usan import dinámico para cargar el MDX:

\begin{lstlisting}[language=JavaScript]
// El import dinámico es resuelto en build time gracias a
// generateStaticParams(), no hay fetch en runtime
const { default: Content } = await import(
  `@/content/inacap/${courseId}/${slug}.mdx`
);
\end{lstlisting}

\begin{mdframed}[style=infobox]
\textbf{Rendimiento}: Al usar \texttt{generateStaticParams()} con \texttt{dynamicParams = false}, Next.js pre-genera HTML para todas las rutas en build time. Las páginas se sirven como archivos HTML estáticos — cero trabajo del servidor en producción.
\end{mdframed}

% =====================================================
\section{Patrones Arquitectónicos Clave}
% =====================================================

\subsection{Patrón: Data + Content Paralelo}

Cada sección del sitio mantiene dos estructuras paralelas:
\begin{itemize}
  \item \textbf{Data} (\texttt{src/data/}): Metadatos, títulos, URLs, fechas, estado de completitud
  \item \textbf{Content} (\texttt{src/content/}): Contenido MDX completo de cada lección
\end{itemize}

Esta separación permite actualizar metadatos (ej: marcar una clase como completada) sin modificar el contenido, y viceversa.

\subsection{Patrón: Filesystem como Fuente de Verdad}

\texttt{getLessonSlugs()} lee el filesystem en lugar de parsear los datos. Esto es más robusto: si existe el archivo MDX, la ruta existe; si no existe, retorna 404. No hay necesidad de mantener sincronizados los datos con el contenido.

\subsection{Patrón: Hook Composition}

Los tres hooks (\texttt{useFilters} → \texttt{useTableSort} → \texttt{usePagination}) se componen en secuencia: los datos fluyen filtrados, luego ordenados, luego paginados. Cada hook es independiente y testeable por separado.

\subsection{Patrón: CSS Classes para Theming}

El sistema de tema no usa React Context ni state management global. Usa clases CSS en el elemento \texttt{<html>} observadas por \texttt{MutationObserver}. Esto permite que cualquier componente (incluso código MDX) pueda reaccionar al tema sin necesidad de prop drilling.

\vspace{1cm}
\begin{center}
\rule{8cm}{0.4pt}\\[0.3cm]
{\small\color{codegray} SaukCode.cl --- Documentación Técnica Vol. III\\
Utilidades, Hooks y Capa de Datos}
\end{center}

\end{document}
